\documentclass[a4paper,11pt]{ltjsarticle}
\usepackage{luatexja}
\usepackage{amsmath,amssymb,amsthm}
\usepackage{longtable,booktabs}
\usepackage{hyperref}
\usepackage[margin=25mm]{geometry}

\newtheorem{theorem}{定理}[section]
\newtheorem{proposition}[theorem]{命題}
\newtheorem{lemma}[theorem]{補題}
\newtheorem{corollary}[theorem]{系}
\newtheorem{definition}[theorem]{定義}
\newtheorem{remark}[theorem]{注意}
\newtheorem{observation}[theorem]{観察}

\providecommand{\tightlist}{\setlength{\itemsep}{0pt}\setlength{\parskip}{0pt}}
\newcounter{none}

\begin{document}

\section{5次元正軸体の配向構造から導かれるスピンの幾何学的分類}\label{ux6b21ux5143ux6b63ux8ef8ux4f53ux306eux914dux5411ux69cbux9020ux304bux3089ux5c0eux304bux308cux308bux30b9ux30d4ux30f3ux306eux5e7eux4f55ux5b66ux7684ux5206ux985e}

\subsection{── 整数幾何学によるスピン値・統計性・力の方向性の統一的導出
──}\label{ux6574ux6570ux5e7eux4f55ux5b66ux306bux3088ux308bux30b9ux30d4ux30f3ux5024ux7d71ux8a08ux6027ux529bux306eux65b9ux5411ux6027ux306eux7d71ux4e00ux7684ux5c0eux51fa}

\textbf{著者:} 木原 範昭 (WF System Co., Ltd.)

\textbf{日付:} 2026年4月

\textbf{DOI:}
\href{https://doi.org/10.5281/zenodo.19630972}{10.5281/zenodo.19630972}

\textbf{前提論文:} -
\href{https://doi.org/10.5281/zenodo.19624957}{離散空間における中心投影の全球被覆と5次元背景空間の整数論的必然性}
-
\href{https://doi.org/10.5281/zenodo.19601592}{万物の理論が満たすべき構造要件の定式化}
-
\href{https://doi.org/10.5281/zenodo.19604965}{標準模型の19個の任意パラメータの分類と構造分析}

\begin{center}\rule{0.5\linewidth}{0.5pt}\end{center}

\subsection{0. 本稿の目的}\label{ux672cux7a3fux306eux76eeux7684}

前稿 {[}W5{]}
において、5次元正軸体（5次元正八面体、5-orthoplex）が離散空間の稠密充填の単位体積を構成し、中心投影により4次元主観空間では超直方体として射影されることを示した。また、体積の正負（配向）が5次元の行列式符号から自然に定義されることを命題7.1で述べた。

本稿は、この配向構造をさらに分析し、以下を導出する。

\begin{enumerate}
\def\labelenumi{\arabic{enumi}.}
\tightlist
\item
  5次元正軸体の配向自由度から、表現可能なスピン量子数が
  \(s = 0, \frac{1}{2}, 1, \frac{3}{2}, 2, \frac{5}{2}\)
  の6種に限定されること
\item
  半整数スピン（フェルミオン）と整数スピン（ボソン）の区別が、頂点状態と辺中点状態の幾何学的差異として導出されること
\item
  体積符号の偶奇構造から、力の方向性（引力のみ／引力斥力両方）が決定されること
\item
  3つの空間軸の置換対称性 \(S_3\)
  から、粒子の世代構造が自然に出現すること
\item
  スピン0の符号非対称性（正10・負9）が、物質／反物質非対称の幾何学的起源を示唆すること
\end{enumerate}

本稿の全ての結果は、離散幾何学と整数論の命題として述べられ、物理的解釈には踏み込まない。ただし、既知の素粒子物理学との対応関係を注記として示す。

\begin{center}\rule{0.5\linewidth}{0.5pt}\end{center}

\subsection{1.
5次元正軸体の基本構造}\label{ux6b21ux5143ux6b63ux8ef8ux4f53ux306eux57faux672cux69cbux9020}

\subsubsection{1.1
定義と要素数}\label{ux5b9aux7fa9ux3068ux8981ux7d20ux6570}

5次元正軸体は、5次元空間 \(\mathbb{R}^5\) において、各座標軸上の
\(\pm e_i\)（\(i = 1, 2, 3, 4, 5\)）を頂点とする正多胞体である。一辺の長さを1とする（頂点間距離
\(\sqrt{2}\) を \(\ell\) とし、\(\ell^2 = 2\) は整数）。

\(k\) 次元面の個数は \(2^{k+1}\binom{5}{k+1}\) で与えられる。

{\def\LTcaptype{none} % do not increment counter
\begin{longtable}[]{@{}cccc@{}}
\toprule\noalign{}
\(k\) & 要素 & 個数 & 計算 \\
\midrule\noalign{}
\endhead
\bottomrule\noalign{}
\endlastfoot
0 & 頂点 & 10 & \(2^1 \binom{5}{1}\) \\
1 & 辺 & 40 & \(2^2 \binom{5}{2}\) \\
2 & 三角形面 & 80 & \(2^3 \binom{5}{3}\) \\
3 & 正四面体胞 & 80 & \(2^4 \binom{5}{4}\) \\
4 & 正5胞体 & 32 & \(2^5 \binom{5}{5}\) \\
\end{longtable}
}

\subsubsection{1.2
計量の整数性}\label{ux8a08ux91cfux306eux6574ux6570ux6027}

一辺を1として各次元の計量を求める。

{\def\LTcaptype{none} % do not increment counter
\begin{longtable}[]{@{}cccc@{}}
\toprule\noalign{}
次元 & 計量 & 1要素あたりの値 & 有理数か \\
\midrule\noalign{}
\endhead
\bottomrule\noalign{}
\endlastfoot
0 & 個数 & 1 & 有理数 \\
1 & 辺長 & 1 & 有理数 \\
2 & 面積（正三角形） & \(\frac{\sqrt{3}}{4}\) & \textbf{無理数} \\
3 & 体積（正四面体） & \(\frac{\sqrt{2}}{12}\) & \textbf{無理数} \\
\end{longtable}
}

0次元（個数）と1次元（辺長）のみが有理数であり、2次元以上の計量は必然的に無理数を含む。このことは、整数幾何学における辺（1次元）の特権的地位を示す。

\subsubsection{1.3
5次元固有の数値的一致}\label{ux6b21ux5143ux56faux6709ux306eux6570ux5024ux7684ux4e00ux81f4}

2次元面（三角形）の個数80と3次元胞（正四面体）の個数80が一致する。これは
\(2^3 \binom{5}{3} = 2^4 \binom{5}{4}\)
による5次元固有の性質であり、面の配向とキラリティが1対1に対応する構造を示唆する。

\begin{center}\rule{0.5\linewidth}{0.5pt}\end{center}

\subsection{2.
配向自由度と独立な向き付き部分空間}\label{ux914dux5411ux81eaux7531ux5ea6ux3068ux72ecux7acbux306aux5411ux304dux4ed8ux304dux90e8ux5206ux7a7aux9593}

\subsubsection{\texorpdfstring{2.1
\(k\)-ベクトルの独立成分数}{2.1 k-ベクトルの独立成分数}}\label{k-ux30d9ux30afux30c8ux30ebux306eux72ecux7acbux6210ux5206ux6570}

5次元空間における独立な向き付き \(k\) 次元部分空間の数は
\(\binom{5}{k}\) で与えられる。

{\def\LTcaptype{none} % do not increment counter
\begin{longtable}[]{@{}cccc@{}}
\toprule\noalign{}
\(k\) & \(\binom{5}{k}\) & 幾何学的意味 & 符号の意味 \\
\midrule\noalign{}
\endhead
\bottomrule\noalign{}
\endlastfoot
0 & 1 & スカラー（存在） & ある／ない \\
1 & 5 & ベクトル（辺の向き） & \(\pm\) 方向 \\
2 & 10 & 二重ベクトル（面の表裏） & \(\pm\) 回転向き \\
3 & 10 & 三重ベクトル（体積の手性） & \(\pm\) キラリティ \\
4 & 5 & 四重ベクトル（超体積の向き） & \(\pm\) 配向 \\
5 & 1 & 擬スカラー（全体配向） & \(\pm\) 全体手性 \\
\end{longtable}
}

合計: \(\sum_{k=0}^{5} \binom{5}{k} = 2^5 = 32\)

この32は、5次元正軸体の4次元胞体（正5胞体）の個数と一致する。

\subsubsection{2.2
パスカルの三角形との関係}\label{ux30d1ux30b9ux30abux30ebux306eux4e09ux89d2ux5f62ux3068ux306eux95a2ux4fc2}

\[1, \; 5, \; 10, \; 10, \; 5, \; 1\]

この対称的な列は、\(k\) と \((5-k)\) のホッジ双対性を反映する。特に
\(\binom{5}{2} = \binom{5}{3} = 10\)
の一致は、2次元の面の表裏と3次元のキラリティが1対1に対応することを意味する。

\begin{center}\rule{0.5\linewidth}{0.5pt}\end{center}

\subsection{3.
スピン値の幾何学的導出}\label{ux30b9ux30d4ux30f3ux5024ux306eux5e7eux4f55ux5b66ux7684ux5c0eux51fa}

\subsubsection{3.1
軸の物理的役割による分類}\label{ux8ef8ux306eux7269ux7406ux7684ux5f79ux5272ux306bux3088ux308bux5206ux985e}

5本の座標軸を対称性に基づいて分類する。

\begin{itemize}
\tightlist
\item
  \textbf{空間軸}: \(x_1, x_2, x_3\)（3本、置換群 \(S_3\) のもとで対称）
\item
  \textbf{第4軸}: \(x_4\)（空間軸とは非対称、符号反転の第1候補）
\item
  \textbf{第5軸}: \(x_5\)（射影の追加次元、符号反転の第2候補）
\end{itemize}

空間軸3本は互いに交換可能であるが、第4軸・第5軸はそれぞれ独立な役割を持つ。

\subsubsection{3.2
符号反転とスピンの対応}\label{ux7b26ux53f7ux53cdux8ee2ux3068ux30b9ux30d4ux30f3ux306eux5bfeux5fdc}

結合行列 \(C\) が \(n\)
本の非空間軸に符号反転を導入するとき、ボソンスピンは \(s = n\) となる。

{\def\LTcaptype{none} % do not increment counter
\begin{longtable}[]{@{}cccc@{}}
\toprule\noalign{}
符号反転数 \(n\) & \(C\) の構造 & ボソンスピン & フェルミオンスピン \\
\midrule\noalign{}
\endhead
\bottomrule\noalign{}
\endlastfoot
0 & \(\mathrm{diag}(+,+,+,+,+)\) & 0 & 1/2 \\
1 & \(\mathrm{diag}(+,+,+,-,+)\) & 1 & 3/2 \\
2 & \(\mathrm{diag}(+,+,+,-,-)\) & 2 & 5/2 \\
\end{longtable}
}

非空間軸は2本（\(x_4, x_5\)）しかないため、符号反転は最大2回である。

\textbf{命題 3.1.} 5次元正軸体の配向構造から表現可能なスピン量子数は

\[s \in \left\{ 0, \; \frac{1}{2}, \; 1, \; \frac{3}{2}, \; 2, \; \frac{5}{2} \right\}\]

の6種に限定される。

\emph{証明.} 空間軸3本は \(S_3\)
対称性により独立な符号反転を与えない。独立に符号反転できる軸は第4軸と第5軸の2本のみである。ボソンスピンは符号反転数に等しいので
\(s = 0, 1, 2\) の3種。フェルミオンスピンは \(s + \frac{1}{2}\)
であるから \(s = \frac{1}{2}, \frac{3}{2}, \frac{5}{2}\)
の3種。合計6種。\(\square\)

\textbf{注記.} 実験で確認されている基本粒子のスピンは
\(0, \frac{1}{2}, 1\) のみであり、\(\frac{3}{2}, 2\)
は理論的予言（グラビティーノ、グラビトン）である。6種すべてが既知の理論的枠組みの範囲内にある。

\subsubsection{3.3
スピン状態数}\label{ux30b9ux30d4ux30f3ux72b6ux614bux6570}

スピン \(s\) の粒子は \(2s + 1\) 個の \(s_z\) 状態を持つ。

{\def\LTcaptype{none} % do not increment counter
\begin{longtable}[]{@{}cc@{}}
\toprule\noalign{}
\(s\) & 状態数 \(2s+1\) \\
\midrule\noalign{}
\endhead
\bottomrule\noalign{}
\endlastfoot
0 & 1 \\
1/2 & 2 \\
1 & 3 \\
3/2 & 4 \\
2 & 5 \\
5/2 & 6 \\
\textbf{合計} & \textbf{21} \\
\end{longtable}
}

\[\sum_{s} (2s+1) = 1 + 2 + 3 + 4 + 5 + 6 = 21 = T(6)\]

21は第6三角数であり、\(\binom{7}{2} = 21\) でもある。

\begin{center}\rule{0.5\linewidth}{0.5pt}\end{center}

\subsection{4.
半整数スピンの幾何学的起源}\label{ux534aux6574ux6570ux30b9ux30d4ux30f3ux306eux5e7eux4f55ux5b66ux7684ux8d77ux6e90}

\subsubsection{4.1
辺の中点と半整数座標}\label{ux8fbaux306eux4e2dux70b9ux3068ux534aux6574ux6570ux5ea7ux6a19}

5次元正軸体の頂点は整数格子上にある（各軸の \(\pm 1\) と \(0\)
の組み合わせ）。

隣接する2頂点、例えば \(e_1 = (1,0,0,0,0)\) と \(e_2 = (0,1,0,0,0)\)
を結ぶ辺の中点は

\[m_{12} = \left(\frac{1}{2}, \frac{1}{2}, 0, 0, 0\right)\]

であり、\textbf{半整数座標}を持つ。

\textbf{命題 4.1.}
整数格子上に定義された5次元正軸体において、辺の中点は半整数格子上に位置する。この半整数格子点が、半整数スピン（フェルミオン）状態の幾何学的居場所である。

\subsubsection{4.2
720°回転の幾何学的意味}\label{ux56deux8ee2ux306eux5e7eux4f55ux5b66ux7684ux610fux5473}

頂点上の状態（ボソン）は、格子上の1周（360°回転）で元の頂点に戻る。

辺の中点上の状態（フェルミオン）は、1周では隣接する中点に移動し、2周（720°回転）で元の中点に戻る。

これはスピン \(\frac{1}{2}\) の粒子が360°回転で位相 \((-1)\)
を獲得し、720°回転で元に戻る性質の幾何学的表現である。

\subsubsection{4.3
パウリ排他原理の幾何学的導出}\label{ux30d1ux30a6ux30eaux6392ux4ed6ux539fux7406ux306eux5e7eux4f55ux5b66ux7684ux5c0eux51fa}

\textbf{命題 4.2.} 半整数スピン状態（辺の中点）では、\(s_z = 0\)
が存在しない。したがって体積の符号は常に \(+\) または \(-\)
のいずれかに確定する。

\emph{証明.} スピン \(s\) の \(z\) 成分は
\(s_z = -s, -s+1, \ldots, s-1, s\) をとる。\(s\)
が半整数のとき、\(s_z = 0\) はこの列に含まれない。\(\square\)

\textbf{系 4.3.}
体積符号が常に確定する状態では、同一の符号・同一の量子数を持つ2つの状態が同一の空間点に共存できない。これはパウリの排他原理に相当する。

対照的に、整数スピン（ボソン）では \(s_z = 0\) が存在し、体積符号が
\(\pm\) 両方を許容するため、同一状態への複数占有が可能である。

\begin{center}\rule{0.5\linewidth}{0.5pt}\end{center}

\subsection{5.
体積符号と力の方向性}\label{ux4f53ux7a4dux7b26ux53f7ux3068ux529bux306eux65b9ux5411ux6027}

\subsubsection{5.1
全スピン状態の体積符号}\label{ux5168ux30b9ux30d4ux30f3ux72b6ux614bux306eux4f53ux7a4dux7b26ux53f7}

各スピン状態 \((s, s_z)\) について、体積の正負を分類する。

\textbf{ボソン状態（整数スピン）:}

{\def\LTcaptype{none} % do not increment counter
\begin{longtable}[]{@{}ccc@{}}
\toprule\noalign{}
\(s\) & \(s_z\) & 体積符号 \\
\midrule\noalign{}
\endhead
\bottomrule\noalign{}
\endlastfoot
0 & 0 & \(+\) \\
1 & \(+1\) & \(+\) \\
1 & \(0\) & \(\pm\) \\
1 & \(-1\) & \(-\) \\
2 & \(+2\) & \(+\) \\
2 & \(+1\) & \(+\) \\
2 & \(0\) & \(\pm\) \\
2 & \(-1\) & \(-\) \\
2 & \(-2\) & \(-\) \\
\end{longtable}
}

\textbf{フェルミオン状態（半整数スピン）:}

{\def\LTcaptype{none} % do not increment counter
\begin{longtable}[]{@{}ccc@{}}
\toprule\noalign{}
\(s\) & \(s_z\) & 体積符号 \\
\midrule\noalign{}
\endhead
\bottomrule\noalign{}
\endlastfoot
1/2 & \(+1/2\) & \(+\) \\
1/2 & \(-1/2\) & \(-\) \\
3/2 & \(+3/2\) & \(+\) \\
3/2 & \(+1/2\) & \(+\) \\
3/2 & \(-1/2\) & \(-\) \\
3/2 & \(-3/2\) & \(-\) \\
5/2 & \(+5/2\) & \(+\) \\
5/2 & \(+3/2\) & \(+\) \\
5/2 & \(+1/2\) & \(+\) \\
5/2 & \(-1/2\) & \(-\) \\
5/2 & \(-3/2\) & \(-\) \\
5/2 & \(-5/2\) & \(-\) \\
\end{longtable}
}

\subsubsection{5.2
符号反転の偶奇と力の方向性}\label{ux7b26ux53f7ux53cdux8ee2ux306eux5076ux5947ux3068ux529bux306eux65b9ux5411ux6027}

\textbf{命題 5.1.} 符号反転数 \(n\)
が偶数のボソン（\(s = 0, 2\)）は、体積符号の積が \((-1)^n = +1\)
となり、力（もし媒介するなら）は一方向的（引力のみ）である。符号反転数が奇数のボソン（\(s = 1\)）は、積が
\((-1)^n = -1\) となり、力は双方向的（引力・斥力の両方）である。

{\def\LTcaptype{none} % do not increment counter
\begin{longtable}[]{@{}cccc@{}}
\toprule\noalign{}
符号反転数 & 積の符号 & ボソンスピン & 力の方向性 \\
\midrule\noalign{}
\endhead
\bottomrule\noalign{}
\endlastfoot
0（偶数） & \((+)(+) = +\) & 0 & 方向なし（スカラー場） \\
1（奇数） & \((+)(-) = -\) & 1 & \(\pm\)（引力・斥力） \\
2（偶数） & \((-)(-)= +\) & 2 & \(+\) のみ（引力のみ） \\
\end{longtable}
}

\textbf{注記.}
電磁力（スピン1媒介）が引力と斥力の両方を持ち、重力（スピン2媒介）が引力のみであることは、実験事実と一致する。

\subsubsection{5.3
スピン0とスピン2の符号の質的差異}\label{ux30b9ux30d4ux30f30ux3068ux30b9ux30d4ux30f32ux306eux7b26ux53f7ux306eux8ceaux7684ux5deeux7570}

スピン0（符号反転0回）とスピン2（符号反転2回）は、ともに体積積が \(+\)
であるが、その由来は本質的に異なる。

{\def\LTcaptype{none} % do not increment counter
\begin{longtable}[]{@{}ccc@{}}
\toprule\noalign{}
& スピン0 & スピン2 \\
\midrule\noalign{}
\endhead
\bottomrule\noalign{}
\endlastfoot
体積積 & \(+\) & \(+\) \\
由来 & \((+)(+)\)：反転なし & \((-)(-)\)：二重反転 \\
幾何学的性質 & 純粋スカラー & 二重反転による復帰 \\
\end{longtable}
}

\textbf{命題 5.2.} \((+)(+)\) 型と \((-)(-)\)
型は、ともに体積積が正であるが、代数的構造が異なる。\((+)(+)\)
型は恒等変換であり、\((-)(-)\) 型は反転の自己逆元である。

\textbf{注記.}
素粒子物理学において、ヒッグス粒子（スピン0）は質量を与える源であり自身も質量を持つが、グラビトン（スピン2）は質量を持たない。この差異は
\((+)(+)\) 型と \((-)(-)\) 型の代数的差異に対応する。

\begin{center}\rule{0.5\linewidth}{0.5pt}\end{center}

\subsection{6.
体積符号の非対称性}\label{ux4f53ux7a4dux7b26ux53f7ux306eux975eux5bfeux79f0ux6027}

\subsubsection{6.1
正負の状態数}\label{ux6b63ux8ca0ux306eux72b6ux614bux6570}

21個の全スピン状態を体積符号で分類する。

{\def\LTcaptype{none} % do not increment counter
\begin{longtable}[]{@{}cc@{}}
\toprule\noalign{}
分類 & 状態数 \\
\midrule\noalign{}
\endhead
\bottomrule\noalign{}
\endlastfoot
常に \(+\) & 10 \\
常に \(-\) & 9 \\
\(\pm\) 両方 & 2 \\
\textbf{合計} & \textbf{21} \\
\end{longtable}
}

\textbf{命題 6.1.}
5次元正軸体の配向構造において、体積が常に正の状態数（10）は、常に負の状態数（9）より1だけ多い。この差はスピン0状態が
\(+\) のみを持つことに起因する。

\subsubsection{6.2
10と5次元正軸体の頂点数}\label{ux30685ux6b21ux5143ux6b63ux8ef8ux4f53ux306eux9802ux70b9ux6570}

体積が常に正の状態数 \textbf{10} は、5次元正軸体の頂点数と一致する。

\textbf{注記.}
観測宇宙において物質が反物質よりわずかに多い（バリオン非対称性）ことは、宇宙論における未解決問題の一つである。本稿の結果は、この非対称性の起源がスピン0（スカラー場）の符号構造にある可能性を示唆する。

\begin{center}\rule{0.5\linewidth}{0.5pt}\end{center}

\subsection{7.
辺の型の分類と粒子の階層構造}\label{ux8fbaux306eux578bux306eux5206ux985eux3068ux7c92ux5b50ux306eux968eux5c64ux69cbux9020}

\subsubsection{\texorpdfstring{7.1 辺の型と \(S_3\)
対称性}{7.1 辺の型と S\_3 対称性}}\label{ux8fbaux306eux578bux3068-s_3-ux5bfeux79f0ux6027}

5次元正軸体の辺は2つの頂点（異なる軸上）を結ぶ。5軸から2軸を選ぶ組み合わせは
\(\binom{5}{2} = 10\) 通りである。

空間軸3本の置換対称性 \(S_3\)
のもとで、10種の辺は以下の4群に分類される。

{\def\LTcaptype{none} % do not increment counter
\begin{longtable}[]{@{}cccc@{}}
\toprule\noalign{}
辺の型 & 軸の組 & 個数 & \(S_3\) 軌道 \\
\midrule\noalign{}
\endhead
\bottomrule\noalign{}
\endlastfoot
空間-空間 & \((x_1, x_2), (x_2, x_3), (x_1, x_3)\) & 3 & 1軌道 \\
空間-第4軸 & \((x_1, x_4), (x_2, x_4), (x_3, x_4)\) & 3 & 1軌道 \\
空間-第5軸 & \((x_1, x_5), (x_2, x_5), (x_3, x_5)\) & 3 & 1軌道 \\
第4軸-第5軸 & \((x_4, x_5)\) & 1 & 1軌道 \\
\end{longtable}
}

\[10 = 3 + 3 + 3 + 1\]

\subsubsection{7.2
3の繰り返しの意味}\label{ux306eux7e70ux308aux8fd4ux3057ux306eux610fux5473}

\(S_3\)
対称性は、3つの空間軸の交換に対する不変性である。各群の中の3つの辺の型は、空間軸の番号
\(i = 1, 2, 3\) のみが異なり、幾何学的に等価である。

\textbf{命題 7.1.} 空間軸の \(S_3\)
対称性のもとで、フェルミオンの辺の型は3つずつのグループを3組形成する。

\textbf{注記.} 素粒子物理学における3世代構造（\(e, \mu, \tau\)
等）は未解決問題の一つである。本稿の結果は、この3世代が3つの空間軸の置換対称性に幾何学的起源を持つ可能性を示唆する。

\subsubsection{7.3
辺の総数との対応}\label{ux8fbaux306eux7dcfux6570ux3068ux306eux5bfeux5fdc}

5次元正軸体の辺の総数は40本である。これは以下のように分解される。

\[40 = 10 \times 2 \times 2\]

{\def\LTcaptype{none} % do not increment counter
\begin{longtable}[]{@{}cc@{}}
\toprule\noalign{}
因子 & 幾何学的意味 \\
\midrule\noalign{}
\endhead
\bottomrule\noalign{}
\endlastfoot
10 & 辺の型（\(\binom{5}{2}\)） \\
\(\times 2\) & 辺の方向（\(\pm\)） \\
\(\times 2\) & 半球の選択（北／南） \\
\end{longtable}
}

\begin{center}\rule{0.5\linewidth}{0.5pt}\end{center}

\subsection{8.
先行研究との関係}\label{ux5148ux884cux7814ux7a76ux3068ux306eux95a2ux4fc2}

\subsubsection{8.1
直接的先行研究}\label{ux76f4ux63a5ux7684ux5148ux884cux7814ux7a76}

5次元正軸体の配向構造からスピン値を導出するアプローチは、著者の知る限り先行研究が存在しない。

\subsubsection{8.2
関連する研究分野}\label{ux95a2ux9023ux3059ux308bux7814ux7a76ux5206ux91ce}

以下の研究が、本稿の個々の側面と関連する。

\textbf{スピンと統計の位相的導出:} - Finkelstein, D. \& Rubinstein, J.
(1968). ``Connection between Spin, Statistics, and Kinks.'' \emph{J.
Math. Phys.} 9, 1762.

\textbf{幾何学的物質モデル:} - Atiyah, M.F., Manton, N.S. \& Schroers,
B.J. (2012). ``Geometric Models of Matter.'' \emph{Proc. R. Soc. A} 468,
1252. arXiv: 1108.5151.

\textbf{クリフォード代数と標準模型:} - Furey, C. (2016). ``Standard
Model Physics from an Algebra?'' arXiv: 1611.09182.

\textbf{5次元統一理論:} - Kaluza, Th. (1921). ``Zum Unitätsproblem der
Physik.'' \emph{Sitzungsber. Preuss. Akad. Wiss. Berlin} 966-972.

\subsubsection{8.3 数値的対応}\label{ux6570ux5024ux7684ux5bfeux5fdc}

本稿で現れる数値と既知の数学的構造の対応を記す。

{\def\LTcaptype{none} % do not increment counter
\begin{longtable}[]{@{}cc@{}}
\toprule\noalign{}
本稿の数値 & 既知の構造 \\
\midrule\noalign{}
\endhead
\bottomrule\noalign{}
\endlastfoot
21（全スピン状態数） & \(\dim \mathrm{SO}(7) = 21\) \\
32（4次元胞体数） & \(\dim \mathrm{Cl}(5) = 2^5 = 32\) \\
10（頂点数） & \(\dim \mathrm{SO}(5) = 10\) \\
\end{longtable}
}

これらの一致が偶然か本質的かの判断は、今後の課題とする。

\begin{center}\rule{0.5\linewidth}{0.5pt}\end{center}

\subsection{9.
未解決問題と今後の課題}\label{ux672aux89e3ux6c7aux554fux984cux3068ux4ecaux5f8cux306eux8ab2ux984c}

本稿の結果は離散幾何学の命題として完結するが、以下の問題が未解決である。

\subsubsection{9.1
定量的質量比の導出}\label{ux5b9aux91cfux7684ux8ceaux91cfux6bd4ux306eux5c0eux51fa}

3つの空間軸の \(S_3\)
対称性が、中心投影の投影軸選択によって破れることで、3世代の質量階層が生じる可能性を示唆した（第7節）。しかし、質量比の具体的な値は導出されていない。

\textbf{課題:} 投影軸と各空間軸のなす角度から、質量比
\(m_e : m_\mu : m_\tau\) 等を定量的に導出できるか。

\subsubsection{\texorpdfstring{9.2 ゲージ群
\(\mathrm{SU}(3) \times \mathrm{SU}(2) \times \mathrm{U}(1)\)
の導出}{9.2 ゲージ群 \textbackslash mathrm\{SU\}(3) \textbackslash times \textbackslash mathrm\{SU\}(2) \textbackslash times \textbackslash mathrm\{U\}(1) の導出}}\label{ux30b2ux30fcux30b8ux7fa4-mathrmsu3-times-mathrmsu2-times-mathrmu1-ux306eux5c0eux51fa}

空間軸の置換群 \(S_3\) はワイル群として \(\mathrm{SU}(3)\)
と関連するが、標準模型のゲージ群全体が5次元正軸体の対称性から導出されるかは未証明である。

\textbf{課題:} \(S_3 \to \mathrm{SU}(3)\), 辺の型の構造
\(\to \mathrm{SU}(2) \times \mathrm{U}(1)\) の対応を厳密に構成できるか。

\subsubsection{9.3
辺の総数40と標準模型の状態数30の差}\label{ux8fbaux306eux7dcfux657040ux3068ux6a19ux6e96ux6a21ux578bux306eux72b6ux614bux657030ux306eux5dee}

1世代あたりの基本フェルミオン状態数は標準模型では15（反粒子含めて30）であるのに対し、5次元正軸体の辺の型×方向×半球は40を与える。差の10の物理的意味が未解明である。

\textbf{課題:}
40と30の差は、右巻きニュートリノ等の未発見状態の存在を予言するのか、あるいは数え方の対応に修正が必要か。

\subsubsection{\texorpdfstring{9.4 \(21 = \dim \mathrm{SO}(7)\)
の意味}{9.4 21 = \textbackslash dim \textbackslash mathrm\{SO\}(7) の意味}}\label{dim-mathrmso7-ux306eux610fux5473}

全スピン状態数21が \(\mathrm{SO}(7)\)
の次元と一致することの意味が不明である。

\textbf{課題:} 5次元正軸体のスピン状態空間は、\(\mathrm{SO}(7)\)
の随伴表現またはその部分空間と同型か。

\subsubsection{9.5
結合定数の幾何学的起源}\label{ux7d50ux5408ux5b9aux6570ux306eux5e7eux4f55ux5b66ux7684ux8d77ux6e90}

本稿はスピン値と力の方向性を導出したが、結合の強さ（結合定数）の値は扱っていない。

\textbf{課題:} 電磁結合定数 \(\alpha \approx 1/137\)
等の値に、5次元正軸体の幾何学的量（角度、面積比等）との対応が存在するか。

\begin{center}\rule{0.5\linewidth}{0.5pt}\end{center}

\subsection{10. 結論}\label{ux7d50ux8ad6}

5次元正軸体の配向構造から、以下が純粋に幾何学的に導出された。

\begin{enumerate}
\def\labelenumi{\arabic{enumi}.}
\tightlist
\item
  \textbf{スピン値の完全分類}:
  \(s = 0, \frac{1}{2}, 1, \frac{3}{2}, 2, \frac{5}{2}\)
  の6種が、5次元における非空間軸2本の符号反転と辺中点の半整数構造から必然的に決定される。
\item
  \textbf{フェルミオン／ボソン区別の導出}: \(s_z = 0\)
  状態の存在／非存在が、整数スピンと半整数スピンの統計的性質の差異を幾何学的に与える。
\item
  \textbf{力の方向性}:
  符号反転回数の偶奇が、媒介する力の方向性（一方向／双方向）を決定する。
\item
  \textbf{粒子階層の起源}: \(S_3\) 対称性による辺の型の
  \(3 + 3 + 3 + 1\) 分解が、世代構造の幾何学的候補を提供する。
\item
  \textbf{物質／反物質非対称}:
  正体積状態10と負体積状態9の差1が、スピン0の符号構造に起因する。
\end{enumerate}

これらの結果は、量子力学のスピンに関する公理（スピン統計定理、パウリ排他原理）が、5次元整数幾何学の定理として再導出できる可能性を示している。

\begin{center}\rule{0.5\linewidth}{0.5pt}\end{center}

\subsection{参考文献}\label{ux53c2ux8003ux6587ux732e}

{[}1{]} 木原範昭, ``中心投影による4次元空間の幾何学的定式化,'' 2026.
DOI: 10.5281/zenodo.19427780

{[}2{]} 木原範昭, ``中心投影の幾何学的対称性：多軸モデルの数学的基盤,''
2026. DOI: 10.5281/zenodo.19434932

{[}3{]} 木原範昭, ``万物の理論が満たすべき構造要件の定式化,'' 2026. DOI:
10.5281/zenodo.19601592

{[}4{]} 木原範昭, ``標準模型の19個の任意パラメータの分類と構造分析,''
2026. DOI: 10.5281/zenodo.19604965

{[}5{]} 木原範昭,
``離散空間における中心投影の全球被覆と5次元背景空間の整数論的必然性,''
2026. DOI: 10.5281/zenodo.19624957

{[}6{]} Finkelstein, D. \& Rubinstein, J. ``Connection between Spin,
Statistics, and Kinks.'' \emph{J. Math. Phys.} 9, 1762 (1968).

{[}7{]} Atiyah, M.F., Manton, N.S. \& Schroers, B.J. ``Geometric Models
of Matter.'' \emph{Proc. R. Soc. A} 468, 1252 (2012). arXiv: 1108.5151.

{[}8{]} Furey, C. ``Standard Model Physics from an Algebra?'' arXiv:
1611.09182 (2016).

{[}9{]} Kaluza, Th. ``Zum Unitätsproblem der Physik.''
\emph{Sitzungsber. Preuss. Akad. Wiss. Berlin} 966-972 (1921).

{[}10{]} Coxeter, H.S.M. \emph{Regular Polytopes.} Dover (1973).

\end{document}
